\chapter{Colostomy}

\section*{Overview}
A colostomy brings part of the colon through an opening in the abdominal wall, called a stoma, so stool can pass into an external pouch.

\section*{Why This Test or Procedure Is Ordered}
It may be temporary or permanent and may be needed for colorectal cancer, bowel obstruction, perforation, inflammatory bowel disease, severe trauma, or protection of a distal bowel repair.

\section*{How This Test or Procedure Is Performed}
During abdominal surgery, the colon is divided and the upstream end is brought through the abdominal wall and secured to the skin. The surgeon may create an end or loop colostomy and may later reverse a temporary stoma.

\section*{Risks and Recovery}
Bleeding, infection, stoma retraction or prolapse, skin irritation, blockage, dehydration, and parastomal hernia are possible. Patients learn pouch changes, skin care, diet, hydration, and warning signs from an ostomy nurse.

\section*{References}
\begin{enumerate}
  \item MedlinePlus. Colostomy. U.S. National Library of Medicine; 2025.
  \item Current Medical Diagnosis and Treatment. McGraw-Hill; 2025.
\end{enumerate}
\clearpage
